\documentclass[11pt,a4paper]{article}

\usepackage[utf8]{inputenc}
\usepackage[T1]{fontenc}
\usepackage[lithuanian]{babel}
\usepackage{lmodern}
\usepackage{geometry}
\usepackage{xcolor}
\usepackage{enumitem}
\usepackage{titlesec}
\usepackage{fancyhdr}
\usepackage{microtype}
\usepackage{ragged2e}

\geometry{
  top=1.8cm,
  bottom=1.8cm,
  left=2.1cm,
  right=2.1cm
}

\definecolor{titleblue}{HTML}{1F3A5F}
\definecolor{ruleblue}{HTML}{6EA8E5}
\definecolor{callout}{HTML}{EEF2F6}

\setlength{\parindent}{0pt}
\setlength{\parskip}{0.45em}
\renewcommand{\familydefault}{\sfdefault}

\titleformat{\section}
  {\large\bfseries\color{titleblue}}
  {}{0pt}{}
\titlespacing*{\section}{0pt}{0.9em}{0.25em}

\titleformat{\subsection}
  {\normalsize\bfseries\color{titleblue}}
  {}{0pt}{}
\titlespacing*{\subsection}{0pt}{0.65em}{0.2em}

\setlist[itemize]{leftmargin=1.3em,itemsep=0.05em,topsep=0.15em,parsep=0pt}
\setlist[enumerate]{leftmargin=1.6em,itemsep=0.08em,topsep=0.15em,parsep=0pt}

\pagestyle{fancy}
\fancyhf{}
\fancyfoot[C]{\thepage}
\renewcommand{\headrulewidth}{0pt}
\renewcommand{\footrulewidth}{0pt}

\newcommand{\calloutbox}[1]{%
  \begin{center}
  \colorbox{callout}{%
    \parbox{0.94\textwidth}{\centering\bfseries #1}%
  }
  \end{center}
}

\begin{document}

\begin{center}
  {\LARGE\bfseries\color{titleblue} Lietuvos mokinių matematikos olimpiados\\[2pt]
  Vilniaus miesto etapas}\par
  \vspace{0.35em}
  {\color{ruleblue}\rule{0.92\textwidth}{0.8pt}}\par
  \vspace{0.45em}
  {\itshape Trumpas gidas pradedančiajam olimpiadininkui}
\end{center}

\section{Kas tai per olimpiada?}

\textbf{Lietuvos mokinių matematikos olimpiados Vilniaus miesto etapas} -- tai pagrindinės Lietuvos \textbf{9--12 klasių mokinių matematikos olimpiados antrasis etapas}.

Visa olimpiada vyksta trimis etapais:
\begin{enumerate}
  \item \textbf{mokyklos etapas};
  \item \textbf{savivaldybės etapas} -- Vilniaus mokiniams tai yra Vilniaus miesto etapas;
  \item \textbf{šalies etapas}.
\end{enumerate}

Taigi Vilniaus miesto etapas nėra atskiras matematikos konkursas. Tai oficialaus kelio į \textbf{Lietuvos mokinių matematikos olimpiados šalies etapą} dalis.

Pradedančiam olimpiadininkui miesto etapas yra pirmas rimtas susitikimas su klasikine olimpiadine matematika: čia svarbu ne tik gauti teisingą atsakymą, bet ir sugebėti \textbf{atrasti sprendimo idėją, ją pagrįsti ir aiškiai užrašyti visą matematinį argumentą}.

\section{Kas gali dalyvauti?}

Pirmiausia mokykloje vyksta \textbf{mokyklos etapas}.

Mokyklos komisija parenka uždavinius, patikrina mokinių darbus ir atrenka geriausiai pasirodžiusius mokinius į Vilniaus miesto etapą.

Pagal dabartinę Vilniaus tvarką iš vienos mokyklos į miesto etapą gali būti siunčiama iki \textbf{3 mokinių iš kiekvienos 9, 10, 11 ir 12 klasės}. Papildomą teisę dalyvauti gali turėti ankstesnių metų miesto ar šalies etapo prizininkai.

Mokinį į olimpiadą registruoja mokykla.

\section{Kiek laiko trunka olimpiada?}

Vilniaus miesto etapo uždaviniams spręsti skiriamos \textbf{4 valandos}.

Tai ilga olimpiada, tačiau uždavinių nėra labai daug. Didžioji laiko dalis skirta ne skaičiavimams, o \textbf{mąstymui, idėjos paieškai ir sprendimo užrašymui}.

\section{Kokia užduočių struktūra?}

Pastarųjų metų savivaldybės etapo rinkiniuose pateikiami \textbf{5 uždaviniai}.

Užduotys skiriamos dviem grupėms:
\begin{itemize}
  \item \textbf{9--10 klasėms};
  \item \textbf{11--12 klasėms}.
\end{itemize}

Tai reiškia, kad devintokas ir dešimtokas sprendžia tą patį uždavinių rinkinį, o vienuoliktokams ir dvyliktokams pateikiamas kitas rinkinys.

Kiekvienas uždavinys vertinamas \textbf{5 taškais}, todėl maksimalus rezultatas yra \textbf{25 taškai}.

\calloutbox{5 uždaviniai $\rightarrow$ 4 valandos $\rightarrow$ po 5 taškus $\rightarrow$ iš viso 25 taškai}

\section{Kokie būna uždaviniai?}

Tai jau \textbf{olimpiadiniai uždaviniai}, o ne įprasti mokykliniai pratimai.

Gali pasitaikyti:
\begin{itemize}
  \item algebros;
  \item skaičių teorijos;
  \item geometrijos;
  \item kombinatorikos;
  \item loginio ir nestandartinio mąstymo uždavinių.
\end{itemize}

Dažnai mokyklinės formulės yra tik įrankiai. Sunkiausia dalis -- suprasti, \textbf{kokią idėją reikia panaudoti}.

Uždavinys gali atrodyti labai trumpas -- vos keli sakiniai -- tačiau jo sprendimo idėjos galima ieškoti pusvalandį ar ilgiau.

Todėl nereikėtų vertinti uždavinio pagal jo sąlygos ilgį.

\section{Kaip vertinami sprendimai?}

Tai viena svarbiausių miesto olimpiados ypatybių.

\textbf{Vien teisingo atsakymo neužtenka.}

Jeigu uždavinio atsakymas yra, pavyzdžiui, 17, vien parašęs „17“ visų taškų negausi.

Reikia parodyti, \textbf{kodėl atsakymas yra 17}.

Vertinamas visas matematinis sprendimas:
\begin{itemize}
  \item teisinga pagrindinė idėja;
  \item logiški tarpiniai žingsniai;
  \item reikalingų teiginių pagrindimas;
  \item visų būtinų atvejų išnagrinėjimas;
  \item teisinga išvada.
\end{itemize}

Pilnai ir taisyklingai išspręstas uždavinys vertas \textbf{5 taškų}.

Tačiau labai svarbu suprasti, kad egzistuoja ir \textbf{daliniai taškai}.

Jeigu viso uždavinio neišsprendei, bet padarei reikšmingą teisingą sprendimo dalį, už ją gali būti skiriama dalis taškų.

\calloutbox{Jeigu turi teisingą matematinę mintį -- užrašyk ją.}

Tuščias lapas garantuoja 0 taškų. Teisinga lema, pastebėtas svarbus ryšys ar įrodyta dalis uždavinio gali būti verta taškų.

\section{Ar reikia išspręsti visus 5 uždavinius?}

\textbf{Ne.}

Tai labai svarbu suprasti pradedančiam olimpiadininkui.

Miesto matematikos olimpiada nėra kontrolinis darbas, kuriame tikimasi surinkti 80--100 procentų.

Uždaviniai specialiai parenkami taip, kad būtų sunkūs.

Todėl vienas visiškai teisingai išspręstas uždavinys jau reiškia \textbf{5 taškus}.

Du pilnai išspręsti uždaviniai -- \textbf{10 taškų}.

Prie jų dar gali prisidėti daliniai taškai iš kitų uždavinių.

Taigi mokinys, kuris neišsprendė visų penkių uždavinių, nebūtinai pasirodė blogai.

Olimpiadoje daug svarbiau \textbf{kiek matematikos sugebėjai padaryti teisingai}, o ne kiek uždavinių pradėjai.

\section{Ar yra minimali riba?}

\textbf{Ne.}

Čia nėra tokios sistemos kaip Matulionio konkurse, kur pirmiausia reikia įveikti tam tikrą I dalies ribą.

Vilniaus miesto etape \textbf{tikrinami visi tavo pateikti sprendimai}.

Galutinis rezultatas yra visų penkių uždavinių taškų suma -- nuo 0 iki 25 taškų.

Todėl net jeigu vieno uždavinio nepavyko išspręsti pilnai, verta kovoti dėl dalinių taškų.

\section{Kaip nustatomos prizinės vietos?}

Fiksuoto skaičiaus, kuris automatiškai reikštų I, II ar III vietą, nėra.

Negalima iš anksto pasakyti, kad tam tikras taškų skaičius garantuos prizinę vietą.

Viskas priklauso nuo konkrečių metų uždavinių sunkumo ir kitų dalyvių rezultatų.

Jeigu variantas labai sunkus, santykinai nedidelis taškų skaičius gali būti labai stiprus rezultatas.

Todėl olimpiadoje neverta galvoti apie procentus taip, kaip kontroliniame darbe.

\section{Kaip patenkama į šalies etapą?}

Vilniaus miesto etapo darbai patikrinami, sudaromi rezultatai ir atrenkami stipriausi miesto mokiniai.

Vilniaus miestui skiriamas ribotas vietų skaičius į šalies etapą, todėl patekimas į jį jau yra labai rimtas rezultatas.

Dėl šių vietų Vilniuje varžosi daug stiprių olimpiadininkų.

Taip pat gali egzistuoti papildomi patekimo į šalies etapą keliai, pavyzdžiui, ankstesnio šalies etapo prizininkams ar pagal kitus olimpiados nuostatuose numatytus kriterijus.

\section{Kuo galima naudotis?}

Olimpiados metu negalima naudotis:
\begin{itemize}
  \item skaičiuotuvu;
  \item mobiliuoju telefonu;
  \item kompiuteriu ar planšete;
  \item formulynu;
  \item knygomis;
  \item lentelėmis.
\end{itemize}

Uždaviniai sprendžiami \textbf{ranka ant popieriaus}.

\section{Kokia turėtų būti strategija per olimpiadą?}

\subsection{1. Perskaityk visus uždavinius}

Pirmomis minutėmis peržiūrėk visus penkis uždavinius.

Nebūtina jų spręsti eilės tvarka.

Gali būti, kad 1-asis uždavinys tau atrodo visiškai neaiškus, o 4-ajame iš karto pamatai idėją.

Pradėk nuo to, kur \textbf{matai daugiausia galimybių pajudėti}.

\subsection{2. Pirmiausia pasiimk užtikrintus taškus}

Jeigu matai uždavinį, kurio sprendimą beveik supranti, verta pirmiausia jį užbaigti.

Vienas pilnas sprendimas suteikia ne tik taškų, bet ir pasitikėjimo likusiai olimpiados daliai.

\subsection{3. Neužstrik prie vieno uždavinio}

Olimpiadiniam uždaviniui galima skirti daug laiko, tačiau jeigu ilgą laiką visiškai nejudi, verta pereiti prie kito.

Kartais grįžęs po valandos į tą patį uždavinį staiga pamatai tai, ko anksčiau nepastebėjai.

\subsection{4. Užrašyk dalinius rezultatus}

Jeigu pilno sprendimo neradai, bet įrodei svarbų faktą, užrašyk jį.

Pavyzdžiui:
\begin{itemize}
  \item įrodei vieną iš dviejų atvejų;
  \item radai reikalingą lygybę;
  \item įrodei svarbią lemą;
  \item nustatei, kokios reikšmės apskritai galimos.
\end{itemize}

Tai gali būti verta dalinių taškų.

\subsection{5. Palik laiko sprendimui sutvarkyti}

Olimpiadoje vertinama ne tai, ką galvojai, o tai, \textbf{ką parašei}.

Todėl sprendimas turi būti toks, kad jį galėtų suprasti žmogus, kuris nematė tavo mąstymo proceso.

Baigęs uždavinį perskaityk sprendimą ir paklausk savęs:

\begin{center}
\textbf{Ar kiekvienas mano žingsnis pagrįstas?}
\end{center}

\section{Kaip ruoštis?}

Vienas geriausių pasiruošimo būdų -- \textbf{spręsti ankstesnių metų Vilniaus miesto arba savivaldybių etapo uždavinius}.

Pradedančiajam nebūtina iš karto spręsti viso varianto per keturias valandas.

Iš pradžių galima paimti \textbf{vieną uždavinį} ir skirti jam daug laiko.

Svarbiausia išmokti:
\begin{itemize}
  \item nepasiduoti po pirmų 10 minučių;
  \item eksperimentuoti;
  \item nagrinėti paprastesnius atvejus;
  \item pastebėti dėsningumus;
  \item formuluoti tarpinius teiginius;
  \item įrodyti savo teiginius.
\end{itemize}

Tik vėliau verta pereiti prie visos olimpiados imitacijos:

\calloutbox{5 uždaviniai $\rightarrow$ 4 valandos $\rightarrow$ 25 galimi taškai}

\section{Koks turėtų būti pirmasis pradedančiojo tikslas?}

Pirmą kartą dalyvaujant miesto olimpiadoje nereikėtų sau kelti tikslo išspręsti visus penkis uždavinius.

Daug prasmingesnis pirmasis tikslas:

\calloutbox{Per keturias valandas pilnai išspręsti bent vieną olimpiadinį uždavinį ir surinkti dalinių taškų kituose.}

Kai vienas pilnas uždavinys tampa įprastu rezultatu, galima siekti dviejų.

Vėliau -- dviejų pilnų ir dalinių taškų kituose.

Taip palaipsniui artėjama prie kovos dėl miesto prizinių vietų ir patekimo į šalies etapą.

Svarbiausia suprasti, kad olimpiadinė matematika nėra greitas teisingų formulių pritaikymas.

Tai gebėjimas \textbf{ilgai būti su problema, kurios sprendimo iš pradžių nežinai, ieškoti, bandyti, klysti, pastebėti naują idėją ir galiausiai ją paversti griežtu matematiniu argumentu}.

\enlargethispage{3\baselineskip}
\begin{center}
\textbf{Būtent šio gebėjimo Vilniaus miesto matematikos olimpiada ir pradeda rimtai reikalauti.}
\end{center}

\thispagestyle{empty}
\end{document}
